% !TeX program = xelatex
\documentclass[14pt,a4paper]{extreport}

\usepackage{fontspec}
\usepackage{polyglossia}
\setdefaultlanguage{russian}
\setotherlanguage{english}
\setmainfont{Times New Roman}
\setsansfont{Arial}
\setmonofont{Courier New}

\usepackage[
    left=30mm,
    right=10mm,
    top=20mm,
    bottom=20mm
]{geometry}
\usepackage{setspace}
\setstretch{1.18}
\usepackage{indentfirst}
\setlength{\parindent}{1.25cm}
\usepackage{titlesec}
\usepackage{tocloft}
\usepackage{array}
\usepackage{longtable}
\usepackage{booktabs}
\usepackage{multirow}
\usepackage{graphicx}
\usepackage{tikz}
\usepackage{float}
\usepackage{caption}
\usepackage{enumitem}
\usepackage{fancyvrb}
\usepackage{hyperref}
\hypersetup{
    unicode=true,
    colorlinks=false,
    pdfborder={0 0 0}
}

\usetikzlibrary{arrows.meta, positioning, shapes.geometric, shapes.multipart}

\renewcommand{\contentsname}{СОДЕРЖАНИЕ}
\renewcommand{\bibname}{СПИСОК ИСПОЛЬЗОВАННЫХ ИСТОЧНИКОВ}
\renewcommand{\chaptername}{ГЛАВА}
\renewcommand{\thechapter}{\arabic{chapter}}
\renewcommand{\thesection}{\thechapter.\arabic{section}}
\renewcommand{\thesubsection}{\thesection.\arabic{subsection}}
\renewcommand{\thetable}{\thechapter.\arabic{table}}
\renewcommand{\thefigure}{\thechapter.\arabic{figure}}

\titleformat{\chapter}[display]
  {\normalfont\bfseries\centering\MakeUppercase}
  {\chaptername\ \thechapter}
  {1ex}
  {\MakeUppercase}
\titleformat{\section}
  {\normalfont\bfseries}
  {\thesection}
  {1em}
  {}
\titleformat{\subsection}
  {\normalfont\bfseries}
  {\thesubsection}
  {1em}
  {}

\captionsetup{labelsep=endash, justification=centering}
\setlist[itemize]{leftmargin=1.25cm}
\setlist[enumerate]{leftmargin=1.25cm}

\begin{document}

\begin{titlepage}
\thispagestyle{empty}
\begin{center}
МИНИСТЕРСТВО ОБРАЗОВАНИЯ РЕСПУБЛИКИ БЕЛАРУСЬ\\
БЕЛОРУССКИЙ ГОСУДАРСТВЕННЫЙ УНИВЕРСИТЕТ\\
ФАКУЛЬТЕТ ПРИКЛАДНОЙ МАТЕМАТИКИ И ИНФОРМАТИКИ\\
Кафедра многопроцессорных систем и сетей
\end{center}

\vspace{42mm}

\begin{center}
\textbf{РЕАЛИЗАЦИЯ WEB-ПРИЛОЖЕНИЯ}\\
\textbf{``УМНЫЙ КАЛЕНДАРЬ''}
\end{center}

\vspace{16mm}

\begin{center}
Курсовой проект
\end{center}

\vspace{32mm}

\begin{flushright}
\begin{minipage}{0.48\textwidth}
Шибитова Николая Дмитриевича\\
обучающегося 3 курса специальности\\
``Прикладная информатика''

\vspace{8mm}

Научный руководитель:\\
заведующий кафедрой многопроцессорных систем и сетей,\\
доцент, кандидат физико-математических наук\\
Андрушкевич А.Н.
\end{minipage}
\end{flushright}

\vfill

\begin{center}
Минск, 2025
\end{center}
\end{titlepage}

\chapter*{ЗАДАНИЕ ПО КУРСОВОМУ ПРОЕКТУ}
\thispagestyle{empty}

\noindent Белорусский государственный университет\\
Факультет прикладной математики и информатики\\
Кафедра многопроцессорных систем и сетей

\vspace{5mm}

\noindent\textbf{Студент} \underline{Шибитов Николай Дмитриевич}

\vspace{3mm}

\noindent\textbf{1. Тема} ``Реализация web-приложения ``Умный календарь''.

\vspace{3mm}

\noindent\textbf{2. Срок представления курсовой работы (курсового проекта) к защите} \underline{10.12.2025}.

\vspace{3mm}

\noindent\textbf{3. Исходные данные к курсовой работе (курсовому проекту):}

\begin{enumerate}
    \item Next.js Documentation: \url{https://nextjs.org/docs}
    \item FastAPI Documentation: \url{https://fastapi.tiangolo.com/}
    \item SQLAlchemy 2.0 Documentation: \url{https://docs.sqlalchemy.org/en/20/}
    \item OpenAI API Documentation: \url{https://platform.openai.com/docs}
    \item Другие литературные источники и электронные ресурсы по теме курсового проекта.
\end{enumerate}

\noindent\textbf{4. Содержание (структура) курсовой работы (курсового проекта)}

\begin{enumerate}
    \item Обзор предметной области и существующих решений для календарного планирования.
    \item Формулировка требований к web-приложению ``Умный календарь'' и проектирование архитектуры.
    \item Реализация клиентской, серверной и интеллектуальной частей приложения.
\end{enumerate}

\noindent\textbf{5. Индивидуальный график выполнения обучающимся курсовой работы (курсового проекта)}

\begin{longtable}{|p{0.68\textwidth}|p{0.22\textwidth}|}
\hline
\textbf{Содержание задания (наименование структурных элементов, этапов работы)} & \textbf{Сроки выполнения} \\
\hline
Анализ предметной области и существующих решений для календарного планирования & до 16.10.2025 \\
\hline
Проектирование структуры web-приложения, базы данных и REST API & до 30.10.2025 \\
\hline
Разработка серверной части приложения с использованием FastAPI и PostgreSQL & до 13.11.2025 \\
\hline
Разработка клиентской части приложения с использованием Next.js и TypeScript & до 27.11.2025 \\
\hline
Интеграция нейросетевого анализа, тестирование приложения. Подготовка отчета и презентации & до 10.12.2025 \\
\hline
\end{longtable}

\newpage
\pagenumbering{arabic}
\setcounter{page}{2}
\tableofcontents

\chapter*{ПЕРЕЧЕНЬ УСЛОВНЫХ ОБОЗНАЧЕНИЙ}
\addcontentsline{toc}{chapter}{ПЕРЕЧЕНЬ УСЛОВНЫХ ОБОЗНАЧЕНИЙ}

\begin{longtable}{p{0.22\textwidth}p{0.68\textwidth}}
API & программный интерфейс приложения (Application Programming Interface).\\
CRUD & набор базовых операций над данными: создание, чтение, обновление, удаление.\\
DTO & объект передачи данных между слоями приложения.\\
HTTP & протокол передачи гипертекста, используемый для обмена между клиентом и сервером.\\
JWT & JSON Web Token, компактный токен для передачи сведений об аутентифицированном пользователе.\\
LLM & большая языковая модель (Large Language Model).\\
ORM & объектно-реляционное отображение, позволяющее работать с таблицами БД через классы языка программирования.\\
REST & архитектурный стиль построения web API.\\
SQL & язык структурированных запросов к реляционным базам данных.\\
UI & пользовательский интерфейс.\\
UX & пользовательский опыт.\\
\end{longtable}

\chapter*{ВВЕДЕНИЕ}
\addcontentsline{toc}{chapter}{ВВЕДЕНИЕ}

Развитие цифровых технологий привело к тому, что значительная часть учебной, рабочей и личной активности человека стала фиксироваться в электронных сервисах. Пользователь планирует занятия, встречи, консультации, дедлайны, личные дела, хранит краткие заметки, ссылки и напоминания. Чем больше таких действий переносится в цифровую среду, тем заметнее становится проблема рациональной организации времени. Особенно остро она проявляется у студентов, которым приходится совмещать учебное расписание, выполнение лабораторных работ, подготовку к зачетам и экзаменам, участие в проектах, подработку и личные задачи.

На рынке существует большое количество календарных и планировочных сервисов. Крупные экосистемы предлагают Google Calendar и Microsoft Outlook, специализированные сервисы помогают согласовывать встречи, а минималистичные планеры позволяют быстро фиксировать задачи на неделю. Однако эти решения не всегда удобны для пользователя, которому нужен простой и независимый инструмент. Часть сервисов перегружена интерфейсом, другая часть ориентирована только на один сценарий, например согласование времени встречи, третья требует привязки к внешней учетной записи и использует пользовательские данные в рамках коммерческой экосистемы.

Предыдущий этап работы по теме ``Умный календарь'' был посвящен преимущественно анализу предметной области, сравнению существующих решений и формированию концепции приложения. В рамках данной курсовой работы основной акцент переносится на практическую реализацию web-приложения. Поэтому в работе рассматривается не только идея и проектирование системы, но и конкретные программные решения: клиентская часть на Next.js и React, серверная часть на FastAPI, хранение данных в PostgreSQL, аутентификация на основе JWT, описание схем данных через Pydantic и интеграция локальной нейросетевой аналитики через OpenAI-compatible API.

Актуальность темы обусловлена необходимостью разработки доступного web-инструмента, который объединяет функции календаря, задач, заметок и аналитики в одном интерфейсе. Практическая значимость работы заключается в создании приложения, которое может использоваться как основа для персонального планировщика, учебного проекта или дальнейшего развития в сторону более сложной системы управления временем.

Целью курсового проекта является реализация web-приложения ``Умный календарь'', предназначенного для ведения личного расписания, управления событиями и заметками, анализа активности пользователя и формирования интеллектуальных рекомендаций.

Для достижения поставленной цели необходимо решить следующие задачи:
\begin{itemize}
    \item проанализировать предметную область и существующие сервисы календарного планирования;
    \item сформулировать функциональные и нефункциональные требования к приложению;
    \item спроектировать архитектуру системы, структуру базы данных и REST API;
    \item реализовать серверную часть приложения с использованием FastAPI, SQLAlchemy, Pydantic и PostgreSQL;
    \item реализовать клиентскую часть приложения с использованием Next.js, React, TypeScript и Tailwind CSS;
    \item реализовать механизм аутентификации и авторизации пользователей на основе JWT;
    \item разработать модуль аналитики и интеграцию с локальной нейросетью;
    \item провести проверку работоспособности основных сценариев приложения.
\end{itemize}

Объектом исследования является процесс организации личного времени пользователя средствами web-приложения. Предметом исследования являются методы проектирования и программной реализации web-системы календарного планирования с поддержкой заметок, аналитики и нейросетевых рекомендаций.

В работе применяются методы анализа существующих программных решений, структурного и объектно-ориентированного проектирования, моделирования данных, разработки REST API, компонентного проектирования пользовательского интерфейса, а также экспериментальной проверки работоспособности реализованных функций.

Курсовая работа состоит из введения, трех глав, заключения, списка использованных источников и приложений. В первой главе рассматривается предметная область и проводится анализ существующих решений. Во второй главе формулируются требования и проектируется архитектура приложения. В третьей главе описывается практическая реализация клиентской и серверной частей, базы данных, аутентификации, аналитики и нейросетевой интеграции.

\chapter{АНАЛИЗ ПРЕДМЕТНОЙ ОБЛАСТИ И СУЩЕСТВУЮЩИХ РЕШЕНИЙ}

\section{Особенности web-приложений для календарного планирования}

Календарное планирование относится к базовым задачам персональной информационной системы. Пользователь ожидает, что приложение позволит быстро увидеть расписание, добавить событие, перенести задачу, зафиксировать заметку и вернуться к ней позже. Несмотря на кажущуюся простоту предметной области, разработка удобного календарного сервиса требует учета нескольких факторов: временной структуры данных, эргономики интерфейса, надежности хранения, разграничения доступа, поддержки разных устройств и скорости выполнения повседневных операций.

Классический календарный сервис строится вокруг сущности события. Событие имеет дату, название, дополнительное описание, тип, а также необязательные время начала и время окончания. Для учебных сценариев этого набора обычно достаточно: пользователь может добавить лекцию, лабораторную работу, консультацию, дедлайн или личное дело. В более сложных системах появляются повторяющиеся события, приглашения участников, несколько календарей, цвета, статусы занятости, напоминания и внешние интеграции. Однако увеличение числа функций не всегда улучшает пользовательский опыт. Для персонального приложения важнее быстрое выполнение типовых действий, чем максимальное количество настроек.

Отдельным классом задач являются заметки. В отличие от события, заметка не всегда связана с точной датой или временем. Она может фиксировать мысль, краткий список дел, ссылку, напоминание или идею, которую пользователь еще не готов превращать в формализованную задачу. Поэтому эффективное приложение для планирования должно уметь работать как со структурированными сущностями календаря, так и с менее формализованными текстовыми записями.

Современные web-приложения для планирования все чаще объединяют несколько подходов: календарь, список задач, заметки, аналитику активности и интеллектуальные подсказки. Это связано с тем, что пользователю важно не только хранить данные, но и понимать свою загрузку: сколько встреч запланировано, какие дни перегружены, какие задачи повторяются, какие темы чаще всего появляются в заметках. Таким образом, ``умный календарь'' можно рассматривать не только как электронную замену бумажного ежедневника, но и как инструмент анализа личной продуктивности.

\section{Анализ существующих решений}

Для определения требований к разрабатываемому приложению были рассмотрены несколько распространенных сервисов: Google Calendar, Doodle и Tweek. Эти решения представляют разные подходы к планированию: полноценный календарь крупной экосистемы, сервис группового согласования времени и минималистичный недельный планер.

Google Calendar является одним из наиболее известных календарных сервисов. Он предоставляет пользователю календарь с различными представлениями, создание событий, напоминания, совместный доступ, интеграцию с Gmail, Google Meet и другими сервисами Google. Сильной стороной является зрелость продукта и широкая интеграция с экосистемой. Для корпоративного и командного использования это серьезное преимущество.

Однако для пользователя, которому нужен простой личный планировщик, Google Calendar может оказаться избыточным. Форма создания события содержит большое количество параметров, часть которых не требуется в учебном или бытовом сценарии. Кроме того, сервис тесно связан с учетной записью Google и внешней инфраструктурой. Это удобно при активном использовании Google Workspace, но снижает независимость пользователя.

Doodle решает другую задачу: согласование времени встречи между несколькими участниками. Организатор создает набор возможных временных слотов, участники отмечают удобные варианты, после чего выбирается итоговое время. Такой подход эффективен для групповых встреч, но сам сервис не является полноценным личным календарем. Он хорошо закрывает один сценарий, но не заменяет ежедневник, список задач и заметки.

Tweek представляет минималистичный недельный планер. Его интерфейс напоминает бумажную страницу с днями недели, куда пользователь записывает краткие задачи. Достоинство такого подхода состоит в простоте и низком пороге входа. Пользователь быстро понимает, как добавить запись, не сталкивается с большим количеством полей и не тратит время на настройку. Недостатком является ограниченность горизонта планирования и слабая поддержка более сложных сценариев: точного времени, аналитики, отдельных заметок, долговременного хранения и структурированной работы с событиями.

\begin{table}[H]
\centering
\caption{Сравнение существующих решений}
\begin{tabular}{|p{0.22\textwidth}|p{0.32\textwidth}|p{0.34\textwidth}|}
\hline
\textbf{Сервис} & \textbf{Преимущества} & \textbf{Недостатки} \\
\hline
Google Calendar & Развитый календарь, совместный доступ, напоминания, интеграция с Google Meet & Перегруженность интерфейса, привязка к экосистеме Google, избыточность для личного использования \\
\hline
Doodle & Удобное согласование времени встреч, простая логика голосования, понятная роль организатора & Не является личным календарем, часть возможностей доступна по подписке, узкая специализация \\
\hline
Tweek & Минималистичный интерфейс, быстрые заметки по дням недели, ощущение бумажного планера & Ограниченный горизонт планирования, слабая работа со структурированными событиями, недостаток аналитики \\
\hline
\end{tabular}
\end{table}

Проведенный анализ показывает, что ни одно из рассмотренных решений полностью не соответствует целям курсового проекта. Требуется приложение, сочетающее простоту Tweek, базовую календарную функциональность Google Calendar и возможность дальнейшего развития в сторону аналитики и интеллектуальных подсказок. При этом приложение должно сохранять независимость от внешних учетных записей и быть реализовано как самостоятельная web-система.

\section{Выявленные проблемы и требования к новому решению}

На основе анализа существующих сервисов можно выделить несколько основных проблем. Первая проблема состоит в перегруженности интерфейса. В крупных календарных сервисах пользователь сталкивается с большим числом параметров, часть которых не нужна для повседневного планирования. Для учебного и личного сценария более уместна компактная форма создания события: название, дата, время, тип и описание.

Вторая проблема связана с фрагментацией инструментов. Календарь, заметки, список задач и аналитика часто существуют в разных приложениях. Пользователь вынужден переключаться между ними, что повышает вероятность потерять информацию или забыть перенести запись в календарь. Поэтому в разрабатываемом приложении календарь и заметки должны быть частью одного информационного пространства.

Третья проблема касается привязки к внешним экосистемам. В популярных сервисах пользовательские данные часто связаны с крупными платформами. Это удобно для синхронизации, но не всегда желательно для учебного проекта и персонального инструмента. В ``Умном календаре'' реализуется собственная регистрация по электронной почте и паролю, а доступ к данным разграничивается средствами приложения.

Четвертая проблема связана с отсутствием персональной аналитики в простых планировщиках. Пользователь видит список событий, но не всегда понимает общую картину: сколько задач приходится на неделю, какие типы событий преобладают, какие дни перегружены. В рамках курсового проекта эта проблема решается с помощью отдельной страницы аналитики и AI-модуля, который формирует текстовый отчет по событиям и заметкам пользователя.

Таким образом, новое приложение должно удовлетворять следующим общим требованиям:
\begin{itemize}
    \item поддерживать регистрацию и вход пользователей;
    \item хранить данные каждого пользователя изолированно;
    \item предоставлять CRUD-операции для событий календаря;
    \item предоставлять CRUD-операции для заметок;
    \item отображать календарь и список актуальных записей в web-интерфейсе;
    \item рассчитывать статистику активности пользователя;
    \item формировать текстовую AI-рекомендацию на основе событий и заметок;
    \item запускаться в локальном окружении с использованием Docker Compose.
\end{itemize}

\section{Выводы по главе 1}

В первой главе была рассмотрена предметная область календарного планирования и проанализированы существующие решения. Было показано, что крупные календарные сервисы предоставляют широкие возможности, но часто оказываются избыточными для личного использования. Специализированные сервисы хорошо решают отдельные задачи, но не заменяют полноценный персональный планировщик. Минималистичные приложения удобны, однако ограничены по функциональности и аналитике.

На основании проведенного анализа сформулирована идея приложения ``Умный календарь'': создать самостоятельное web-приложение, объединяющее календарь, заметки, аналитику и интеллектуальные рекомендации в одном интерфейсе. Основной акцент в данной курсовой работе делается на практическую реализацию такого приложения.

\chapter{ПРОЕКТИРОВАНИЕ WEB-ПРИЛОЖЕНИЯ ``УМНЫЙ КАЛЕНДАРЬ''}

\section{Назначение и пользовательские сценарии}

Разрабатываемое приложение предназначено для пользователя, которому необходимо вести личное расписание, фиксировать задачи и заметки, а также получать представление о собственной загрузке. Базовым пользователем можно считать студента, так как его расписание сочетает регулярные занятия, разовые консультации, дедлайны по лабораторным работам, подготовку к конференциям и личные дела. Тем не менее архитектура приложения не привязана к учебной предметной области и может применяться для более широких задач персонального планирования.

Основной сценарий начинается с регистрации или входа пользователя. После успешной аутентификации пользователь попадает в основную область приложения, где видит календарь, события выбранного периода и блок текущих задач. Далее он может создать новое событие, отредактировать существующее, удалить устаревшую запись, добавить заметку или перейти на страницу аналитики. Если access-токен истек, клиентская часть автоматически обращается к endpoint обновления токена и продолжает работу без повторного ввода пароля, пока refresh-токен действителен.

Второй важный сценарий связан с заметками. Пользователь может создать заметку без строгого времени выполнения, например ``уточнить тему доклада'' или ``проверить список литературы''. Такая запись хранится отдельно от событий, но может быть связана с датой. Это позволяет сохранить гибкость: заметка не требует полного заполнения формы события, но остается в одном приложении с календарем.

Третий сценарий относится к аналитике. Пользователь открывает страницу аналитики и получает статистику по количеству событий, распределению по типам, активности по дням и числу заметок. Дополнительно он может запросить AI-анализ: сервер собирает события и заметки, формирует промпт и отправляет его локальной языковой модели. Ответ отображается в интерфейсе как краткий отчет с рекомендациями.

\section{Функциональные требования}

Функциональные требования определяют, какие действия система должна выполнять. В рамках курсового проекта реализованы следующие группы функций.

Первая группа функций относится к аутентификации. Пользователь должен иметь возможность зарегистрироваться по электронной почте, имени и паролю. При входе приложение должно проверять пароль, возвращать access- и refresh-токены, а также предоставлять endpoint для получения данных текущего пользователя. Пароль не должен храниться в открытом виде; вместо него сохраняется bcrypt-хэш.

Вторая группа функций связана с событиями календаря. Система должна позволять создавать, просматривать, редактировать и удалять события. Каждое событие содержит название, дату, необязательное время начала и окончания, описание, тип и ссылку. Тип события используется для группировки и аналитики; в проекте предусмотрены значения meeting, task и note.

Третья группа функций относится к заметкам. Заметка содержит заголовок, текст, дату создания и необязательную привязку к календарной дате. Пользователь должен иметь возможность создать, получить, изменить и удалить заметку.

Четвертая группа функций связана с аналитикой. Система должна подсчитывать общее количество событий за период, распределение по типам, ежедневную активность и количество заметок пользователя. Эти данные используются как для визуального отображения статистики, так и для подготовки входного контекста нейросетевого анализа.

Пятая группа функций относится к AI-модулю. Сервер должен уметь собрать события и заметки текущего пользователя, преобразовать их в текстовый контекст, отправить запрос в локальную языковую модель и вернуть полученный отчет клиентской части.

\begin{longtable}{|p{0.18\textwidth}|p{0.72\textwidth}|}
\caption{Функциональные требования к приложению}\\
\hline
\textbf{Код} & \textbf{Описание требования}\\
\hline
FR-1 & Пользователь может зарегистрироваться с указанием имени, электронной почты и пароля.\\
\hline
FR-2 & Пользователь может выполнить вход и получить пару JWT-токенов.\\
\hline
FR-3 & Система позволяет обновить access-токен по refresh-токену.\\
\hline
FR-4 & Пользователь может создать, просмотреть, изменить и удалить событие календаря.\\
\hline
FR-5 & Пользователь может фильтровать события по диапазону дат.\\
\hline
FR-6 & Пользователь может создать, просмотреть, изменить и удалить заметку.\\
\hline
FR-7 & Система отображает статистику активности пользователя.\\
\hline
FR-8 & Система формирует AI-анализ на основе событий и заметок пользователя.\\
\hline
FR-9 & Данные разных пользователей изолированы друг от друга.\\
\hline
\end{longtable}

\section{Нефункциональные требования}

Нефункциональные требования описывают качество системы. Для ``Умного календаря'' наиболее важны удобство интерфейса, надежность хранения данных, безопасность доступа и сопровождаемость кода.

Удобство интерфейса достигается за счет разделения приложения на понятные страницы: вход и регистрация, главный календарь, заметки, аналитика. Визуальный стиль строится на Tailwind CSS, что позволяет быстро задавать единые отступы, цвета, состояния элементов и адаптивное поведение. Пользователь не должен изучать сложную инструкцию: основные действия представлены кнопками, формами и привычными карточками.

Надежность хранения обеспечивается применением PostgreSQL. Реляционная база данных подходит для структурированных сущностей: пользователей, событий, заметок, шаблонов и уведомлений. Между таблицами используются внешние ключи, что помогает поддерживать целостность данных. SQLAlchemy позволяет описывать таблицы в виде Python-классов, а приложение работает с объектами, не формируя SQL-запросы вручную в каждом обработчике.

Безопасность доступа реализуется с помощью JWT. Каждый защищенный endpoint получает текущего пользователя через FastAPI dependency. Это означает, что бизнес-логика не принимает user\_id от клиента как доверенный параметр, а извлекает его из токена. Такой подход снижает риск получения чужих данных при подмене параметров запроса.

Сопровождаемость обеспечивается разделением проекта на backend и frontend, а также внутренним разделением backend на роутеры, модели, схемы, сервисы и утилиты. Клиентская часть использует TypeScript, благодаря чему типы календарных событий, заметок и ответов API фиксируются на этапе разработки.

\begin{longtable}{|p{0.22\textwidth}|p{0.68\textwidth}|}
\caption{Нефункциональные требования}\\
\hline
\textbf{Характеристика} & \textbf{Требование}\\
\hline
Безопасность & Пароли хранятся в виде bcrypt-хэшей; доступ к данным выполняется только по JWT.\\
\hline
Сопровождаемость & Код разделен на независимые модули: роутеры, модели, схемы, сервисы и компоненты интерфейса.\\
\hline
Расширяемость & Архитектура допускает добавление новых типов событий, уведомлений, внешних интеграций и новых AI-функций.\\
\hline
Удобство & Основные действия пользователя выполняются без сложной навигации и перегруженных форм.\\
\hline
Переносимость & Backend и база данных запускаются через Docker Compose.\\
\hline
\end{longtable}

\section{Архитектура приложения}

Приложение реализовано в виде классической клиент-серверной web-системы. Клиентская часть отвечает за отображение интерфейса и взаимодействие с пользователем. Серверная часть предоставляет REST API, выполняет проверку токенов, реализует бизнес-логику и обращается к базе данных. PostgreSQL хранит состояние системы. Отдельным внешним по отношению к приложению компонентом выступает локальная нейросетевая среда Ollama, доступ к которой выполняется через OpenAI-compatible API.

\begin{figure}[H]
\centering
\begin{tikzpicture}[
    node distance=18mm,
    box/.style={rectangle, draw, rounded corners, align=center, minimum width=45mm, minimum height=13mm},
    arrow/.style={-{Latex[length=3mm]}, thick}
]
\node[box] (client) {Frontend\\Next.js, React, TypeScript};
\node[box, below=of client] (api) {Backend\\FastAPI, Pydantic, SQLAlchemy};
\node[box, below left=18mm and 8mm of api] (db) {PostgreSQL\\users, events, notes};
\node[box, below right=18mm and 8mm of api] (ai) {Ollama\\локальная LLM};

\draw[arrow] (client) -- node[right]{HTTP/JSON} (api);
\draw[arrow] (api) -- node[left]{SQLAlchemy ORM} (db);
\draw[arrow] (api) -- node[right]{OpenAI-compatible API} (ai);
\end{tikzpicture}
\caption{Общая архитектура приложения}
\end{figure}

Выбранная архитектура имеет несколько преимуществ. Во-первых, frontend и backend могут разрабатываться независимо. Во-вторых, серверное API потенциально может обслуживать не только web-клиент, но и мобильное приложение. В-третьих, база данных изолирована от клиентской части: пользователь не обращается к ней напрямую, а все операции проходят через слой проверки и бизнес-логики.

\section{Проектирование базы данных}

База данных проектировалась исходя из основных сущностей предметной области: пользователь, событие, заметка, шаблон события и уведомление. Центральной сущностью является пользователь. Все персональные данные, создаваемые в приложении, связаны с пользователем через внешний ключ. Это позволяет реализовать изоляцию данных: пользователь видит только свои события и заметки.

Событие описывает календарную запись. Оно связано с пользователем, имеет дату, название, описание, тип, время начала и окончания. Тип события нужен для аналитики: система может показать, сколько встреч, задач и заметок-событий было создано за выбранный период. Поле link позволяет прикреплять ссылку, например на онлайн-встречу или учебный материал.

Заметка хранит свободный текст. Она также принадлежит пользователю и может иметь необязательную дату. Это отличие от события принципиально: заметка может существовать без точной привязки к расписанию. Такой подход сохраняет гибкость блокнота.

Шаблон события нужен для дальнейшего расширения приложения. Он позволяет хранить повторяемые параметры типовых событий: название, тип, длительность, описание и ссылку. Например, пользователь может создать шаблон ``Лабораторная работа'' или ``Консультация''.

Уведомление связывается с пользователем и событием. В текущей реализации оно хранит сообщение, время отправки и признаки прочтения/отправки. Даже если полноценная рассылка уведомлений не является центральной частью курсового проекта, наличие такой таблицы показывает, что архитектура рассчитана на развитие.

\begin{longtable}{|p{0.2\textwidth}|p{0.2\textwidth}|p{0.48\textwidth}|}
\caption{Основные таблицы базы данных}\\
\hline
\textbf{Таблица} & \textbf{Назначение} & \textbf{Ключевые поля}\\
\hline
users & Хранение учетных записей & id, email, name, password\_hash, avatar\_url, created\_at\\
\hline
events & Хранение календарных событий & id, user\_id, title, description, date, start\_time, end\_time, type, link\\
\hline
notes & Хранение заметок пользователя & id, user\_id, title, text, date, created\_at\\
\hline
event\_templates & Хранение шаблонов событий & id, user\_id, title, type, default\_duration, description\\
\hline
notifications & Хранение уведомлений & id, user\_id, event\_id, message, is\_read, send\_at, sent\\
\hline
\end{longtable}

\section{Проектирование REST API}

REST API является границей между frontend и backend. Клиентская часть не знает внутренней структуры базы данных и не выполняет SQL-запросы. Вместо этого она отправляет HTTP-запросы на серверные endpoint'ы. Ответы передаются в формате JSON, что удобно для TypeScript-клиента и для автоматической валидации Pydantic-схемами.

API разделено на несколько логических групп: аутентификация, события, заметки, шаблоны, уведомления, аналитика и AI. Каждая группа реализуется отдельным роутером FastAPI. Такой подход улучшает читаемость кода и упрощает сопровождение. Например, изменения в логике заметок не требуют просмотра всего файла с серверными обработчиками.

\begin{longtable}{|p{0.24\textwidth}|p{0.18\textwidth}|p{0.46\textwidth}|}
\caption{Ключевые endpoint'ы приложения}\\
\hline
\textbf{Endpoint} & \textbf{Метод} & \textbf{Назначение}\\
\hline
/api/auth/register & POST & Регистрация нового пользователя\\
\hline
/api/auth/login & POST & Вход и получение JWT-токенов\\
\hline
/api/auth/refresh & POST & Обновление access-токена\\
\hline
/api/auth/me & GET & Получение данных текущего пользователя\\
\hline
/api/events & GET & Получение событий пользователя за период\\
\hline
/api/events & POST & Создание события\\
\hline
/api/events/\{id\} & PUT & Обновление события\\
\hline
/api/events/\{id\} & DELETE & Удаление события\\
\hline
/api/notes & GET & Получение заметок пользователя\\
\hline
/api/notes & POST & Создание заметки\\
\hline
/api/analytics & GET & Получение статистики активности\\
\hline
/api/ai/analysis & GET & Получение AI-анализа продуктивности\\
\hline
\end{longtable}

Важной особенностью является то, что большинство endpoint'ов защищено. Текущий пользователь определяется не из параметров URL, а из Bearer-токена. Это делает API устойчивее к ошибкам клиентской части и уменьшает вероятность доступа к чужим данным.

\section{Выбор технологического стека}

Для frontend выбран Next.js 14 на базе React и TypeScript. React обеспечивает компонентный подход: интерфейс делится на переиспользуемые элементы, такие как календарь, карточка события, форма авторизации, виджет AI-анализа. Next.js предоставляет файловую маршрутизацию и удобную структуру приложения. TypeScript позволяет явно описывать типы данных, приходящих с сервера, и уменьшает количество ошибок при изменении API.

Для стилизации выбран Tailwind CSS. Его преимущество состоит в том, что стили задаются непосредственно через utility-классы. В небольшом и среднем web-приложении это ускоряет разработку и помогает сохранять визуальную согласованность без создания большого количества отдельных CSS-файлов.

Для backend выбран FastAPI. Этот фреймворк хорошо подходит для разработки REST API на Python: он поддерживает dependency injection, автоматическую генерацию OpenAPI-документации, Pydantic-валидацию и асинхронные обработчики. FastAPI лаконичен, но при этом достаточно строг для построения реального приложения.

SQLAlchemy используется как ORM. Это позволяет описывать таблицы базы данных Python-классами и работать с объектами вместо ручного формирования SQL-запросов. PostgreSQL выбран как надежная реляционная СУБД, подходящая для хранения структурированных данных приложения.

Docker Compose используется для запуска backend и базы данных в согласованном окружении. Это важно для курсового проекта, так как упрощает воспроизведение результата: преподаватель или другой разработчик может поднять сервер и БД одной командой, не настраивая вручную PostgreSQL.

\section{Проектирование пользовательского интерфейса}

Пользовательский интерфейс проектировался исходя из того, что приложение должно использоваться регулярно, а не только демонстрироваться один раз. Поэтому основной акцент был сделан на быстрый доступ к главным действиям: открыть календарь, создать событие, перейти к заметкам, посмотреть аналитику и получить AI-рекомендацию. Интерфейс не должен превращаться в сложную административную панель, так как предметная область приложения связана с личной организацией времени.

В качестве базовой структуры выбран набор отдельных страниц. Страница авторизации отделена от основного приложения, так как до входа пользователь не должен видеть персональные данные. После входа пользователь попадает в календарную часть, где сосредоточены события и основные действия с расписанием. Раздел заметок вынесен отдельно, поскольку заметки могут не иметь точной даты и являются самостоятельной формой хранения информации. Раздел аналитики служит для обобщения уже накопленных данных.

\begin{longtable}{|p{0.2\textwidth}|p{0.28\textwidth}|p{0.4\textwidth}|}
\caption{Основные экраны приложения}\\
\hline
\textbf{Экран} & \textbf{Назначение} & \textbf{Ключевые элементы}\\
\hline
Авторизация & Регистрация нового пользователя и вход в систему & Форма email/password, переключение режимов входа и регистрации, обработка ошибок\\
\hline
Календарь & Работа с расписанием пользователя & Календарная сетка, список событий, модальное окно события, кнопки создания и редактирования\\
\hline
Заметки & Хранение свободных текстовых записей & Список заметок, форма создания заметки, дата заметки, текстовое содержимое\\
\hline
Аналитика & Просмотр статистики активности & Карточки показателей, распределение по типам, активность по дням, AI-виджет\\
\hline
Профиль/верхняя панель & Навигация и пользовательское состояние & Имя пользователя, переходы между разделами, выход из аккаунта\\
\hline
\end{longtable}

При проектировании интерфейса также учитывались состояния загрузки и ошибки. Для web-приложения, взаимодействующего с API, важно не только показать успешный результат, но и корректно обработать ситуацию, когда сервер недоступен, токен истек или данные не прошли валидацию. Поэтому на уровне клиентского API предусмотрена обработка неуспешных ответов, а на уровне страниц --- вывод понятных сообщений пользователю.

Отдельным решением стало использование модальных окон для создания и редактирования событий. Такой подход позволяет не уводить пользователя с календарной страницы: он видит контекст выбранного дня и одновременно может изменить запись. Для заметок выбран более простой вариант со страницей списка, потому что заметка обычно содержит больше текста и не всегда связана с конкретным временем.

\section{Проектирование обработки ошибок}

Обработка ошибок является важной частью проектирования, так как приложение использует несколько компонентов: frontend, backend, базу данных и локальную нейросеть. Ошибка может возникнуть на любом из этих уровней, поэтому она должна быть локализована и преобразована в понятное сообщение.

На уровне клиента возможны ошибки сетевого взаимодействия, например недоступность backend или неверный адрес API. В такой ситуации браузер может показать сообщение вида \texttt{NetworkError when attempting to fetch resource}. Для уменьшения вероятности таких проблем клиентская часть должна использовать единый модуль обращения к API, где централизованно задается базовый адрес сервера. Это лучше, чем прописывать адреса запросов в каждой странице отдельно.

На уровне backend ошибки делятся на ошибки валидации, ошибки авторизации и ошибки бизнес-логики. Ошибка валидации возникает, если клиент отправил JSON без обязательного поля или с неверным типом данных. FastAPI и Pydantic автоматически формируют ответ с описанием проблемы. Ошибка авторизации возникает при отсутствии токена, истечении срока действия access-токена или попытке использовать refresh-токен вместо access-токена. Ошибка бизнес-логики возникает, например, если пользователь пытается обновить несуществующее событие.

\begin{longtable}{|p{0.23\textwidth}|p{0.28\textwidth}|p{0.37\textwidth}|}
\caption{Типовые ошибки и способы их обработки}\\
\hline
\textbf{Уровень} & \textbf{Пример ошибки} & \textbf{Обработка}\\
\hline
Frontend & Backend недоступен & Показ сообщения о проблеме соединения, повтор запроса после восстановления сервера\\
\hline
Frontend & Истек access-токен & Автоматический запрос \texttt{/api/auth/refresh}, повтор исходного запроса\\
\hline
Backend & Некорректный JSON & Ответ 422 с описанием ошибки валидации Pydantic\\
\hline
Backend & Пользователь не авторизован & Ответ 401, перенаправление клиента на страницу входа\\
\hline
База данных & Нарушение уникальности email & Ответ о том, что пользователь с таким email уже существует\\
\hline
AI-сервис & Ollama не запущена & Понятное сообщение о недоступности локальной нейросети\\
\hline
\end{longtable}

Такое проектирование повышает надежность приложения. Пользователь не должен видеть внутренние stack trace или технические исключения, так как они не помогают решить проблему. Вместо этого приложение должно объяснить, какое действие требуется: повторить вход, запустить локальную нейросеть, проверить соединение с сервером или исправить данные формы.

\section{Трассировка требований к реализации}

Для курсового проекта важно показать связь между сформулированными требованиями и фактически разработанными модулями. Такая трассировка помогает убедиться, что реализация не является набором несвязанных экранов, а соответствует исходной постановке задачи.

\begin{longtable}{|p{0.18\textwidth}|p{0.34\textwidth}|p{0.36\textwidth}|}
\caption{Связь требований с компонентами реализации}\\
\hline
\textbf{Требование} & \textbf{Компонент backend} & \textbf{Компонент frontend}\\
\hline
FR-1, FR-2 & \texttt{routers/auth.py}, \texttt{auth.py}, модель \texttt{User} & Страница \texttt{auth}, формы входа и регистрации\\
\hline
FR-3, FR-4 & \texttt{routers/events.py}, модель \texttt{Event}, схемы \texttt{EventCreate/EventOut} & Главная страница календаря, модальное окно события\\
\hline
FR-5 & \texttt{routers/notes.py}, модель \texttt{Note} & Страница заметок, форма заметки\\
\hline
FR-6 & \texttt{routers/templates.py}, модель \texttt{EventTemplate} & Заложено для последующего интерфейса шаблонов\\
\hline
FR-7 & \texttt{routers/analytics.py}, схема \texttt{AnalyticsOut} & Страница аналитики, карточки статистики\\
\hline
FR-8 & \texttt{routers/ai.py}, \texttt{services/ai\_service.py} & AI-виджет на странице аналитики\\
\hline
FR-9 & Dependency текущего пользователя, фильтрация по \texttt{user\_id} & Хранение токенов и запросы от имени текущего пользователя\\
\hline
\end{longtable}

Из таблицы видно, что каждое ключевое требование имеет отражение как на стороне сервера, так и на стороне клиента. Исключение составляют функции, заложенные архитектурно для будущего развития, например расширенный интерфейс шаблонов и уведомлений. Их наличие в модели данных позволяет развивать приложение без радикального изменения структуры.

\section{Выводы по главе 2}

Во второй главе были определены назначение приложения, пользовательские сценарии, функциональные и нефункциональные требования. Спроектирована клиент-серверная архитектура, структура базы данных и REST API. Обоснован выбор технологического стека: Next.js и React для интерфейса, FastAPI и Pydantic для API, SQLAlchemy и PostgreSQL для хранения данных, Docker Compose для запуска окружения и локальная LLM для интеллектуального анализа.

\chapter{РЕАЛИЗАЦИЯ WEB-ПРИЛОЖЕНИЯ}

\section{Структура проекта}

Проект организован как монорепозиторий, в котором клиентская и серверная части находятся в отдельных каталогах. Такой подход удобен для учебной разработки: вся система хранится в одной папке, но логические границы между frontend и backend остаются явными.

Серверная часть расположена в каталоге \texttt{backend}. В ней находятся точка входа FastAPI-приложения, подключение к базе данных, ORM-модели, Pydantic-схемы, роутеры и сервис AI-анализа. Клиентская часть расположена в каталоге \texttt{frontend}. Она содержит приложение Next.js, страницы, компоненты, типы и вспомогательные функции для обращения к API.

\begin{Verbatim}[fontsize=\small]
Kursovoi_Project/
  backend/
    main.py
    db.py
    models.py
    schemas.py
    auth.py
    routers/
      auth.py
      events.py
      notes.py
      analytics.py
      ai.py
    services/
      ai_service.py
  frontend/
    src/
      app/
        auth/
        analytics/
        notes/
      components/
      lib/
        api.ts
      types/
  docker-compose.yml
\end{Verbatim}

Такое разделение позволяет рассматривать приложение как набор слоев. Верхний слой --- пользовательский интерфейс. Средний слой --- REST API и бизнес-логика. Нижний слой --- база данных. Отдельно расположен сервис локальной нейросети, к которому backend обращается при необходимости.

\section{Реализация серверной части}

Точкой входа серверной части является FastAPI-приложение. В нем регистрируются роутеры, настраивается CORS и создается приложение с описанием. Каждый роутер отвечает за свою предметную область. Например, роутер \texttt{events.py} реализует операции над событиями, \texttt{notes.py} --- операции над заметками, \texttt{analytics.py} --- статистику, \texttt{ai.py} --- обращение к интеллектуальному модулю.

Работа с базой данных организована через SQLAlchemy. В файле \texttt{db.py} создается engine, базовый класс моделей и dependency \texttt{get\_session}. Dependency возвращает сессию базы данных для обработки запроса и закрывает ее после завершения. Это соответствует типичной практике FastAPI: обработчик endpoint'а не создает подключение вручную, а получает готовую сессию через механизм зависимостей.

Модели SQLAlchemy описывают таблицы приложения. Например, модель пользователя содержит email, имя, хэш пароля и дату создания. Модель события содержит дату, время, тип и ссылку на пользователя. Связи \texttt{relationship} позволяют обращаться к событиям пользователя как к объектной коллекции, а каскадное удаление помогает автоматически удалять зависимые записи при удалении пользователя.

\begin{figure}[H]
\centering
\begin{tikzpicture}[
    entity/.style={rectangle split, rectangle split parts=2, draw, align=left, minimum width=38mm},
    rel/.style={-{Latex[length=2.5mm]}, thick}
]
\node[entity] (u) {
    \textbf{User}
    \nodepart{second}
    id\\email\\name\\password\_hash
};
\node[entity, right=18mm of u] (e) {
    \textbf{Event}
    \nodepart{second}
    id\\user\_id\\title\\date\\type
};
\node[entity, below=14mm of e] (n) {
    \textbf{Note}
    \nodepart{second}
    id\\user\_id\\title\\text\\date
};
\node[entity, below=14mm of u] (t) {
    \textbf{EventTemplate}
    \nodepart{second}
    id\\user\_id\\title\\duration
};
\draw[rel] (u) -- node[above]{1:N} (e);
\draw[rel] (u) -- node[below left]{1:N} (n);
\draw[rel] (u) -- node[left]{1:N} (t);
\end{tikzpicture}
\caption{Упрощенная модель данных приложения}
\end{figure}

Серверная часть не смешивает ORM-модели и схемы API. ORM-модель описывает, как данные хранятся в базе, а Pydantic-схема описывает, какие данные приходят от клиента и какие данные возвращаются в ответе. Это важное разделение: внутренняя структура таблицы может содержать поля, которые нельзя отдавать клиенту, например \texttt{password\_hash}. Pydantic-схемы позволяют явно контролировать такие границы.

\section{Использование FastAPI и Pydantic-схем}

FastAPI тесно связан с Pydantic. Когда клиент отправляет JSON-запрос, FastAPI преобразует его в объект Pydantic-схемы и проверяет типы полей. Если данные некорректны, сервер автоматически возвращает ошибку валидации. Это упрощает код обработчиков: внутри функции можно работать уже с проверенными данными.

В проекте используются отдельные схемы для создания, обновления и вывода сущностей. Например, для события определены \texttt{EventCreate}, \texttt{EventUpdate} и \texttt{EventOut}. Схема создания содержит обязательные поля, необходимые для нового события. Схема обновления делает поля необязательными, так как пользователь может изменить только название или только дату. Схема вывода содержит идентификатор, user\_id и дату создания, то есть поля, которые формируются на сервере.

Такой подход решает сразу несколько задач. Во-первых, он документирует API: по схеме видно, какие поля ожидает endpoint. Во-вторых, он защищает сервер от некорректных данных. В-третьих, он помогает frontend-разработке, так как структура ответов становится стабильной и предсказуемой.

Pydantic v2 также используется для преобразования ORM-объектов в ответы API. Параметр \texttt{model\_config = \{"from\_attributes": True\}} позволяет формировать выходную схему на основе атрибутов SQLAlchemy-модели. Благодаря этому endpoint может вернуть ORM-объект, а FastAPI преобразует его в JSON-ответ нужного вида.

Пример логики endpoint'а для создания события можно описать следующим образом: сервер получает объект \texttt{EventCreate}, извлекает текущего пользователя из JWT, создает SQLAlchemy-объект \texttt{Event}, сохраняет его в базе данных, выполняет commit и возвращает объект в формате \texttt{EventOut}. На каждом этапе разделены обязанности: аутентификация отвечает за пользователя, Pydantic --- за входные данные, SQLAlchemy --- за сохранение, FastAPI --- за HTTP-ответ.

\section{Реализация аутентификации и JWT}

Аутентификация реализована по схеме email/password. При регистрации пользователь передает имя, электронную почту и пароль. Сервер проверяет, что пользователь с таким email еще не существует, хэширует пароль с помощью bcrypt и сохраняет новую запись в таблице users. Пароль в открытом виде в базе данных не хранится.

При входе пользователь отправляет email и пароль. Сервер находит пользователя по email, сравнивает переданный пароль с сохраненным хэшем и при успешной проверке создает пару токенов: access-токен и refresh-токен. Access-токен имеет короткое время жизни и используется для доступа к защищенным endpoint'ам. Refresh-токен имеет более длительное время жизни и используется для получения нового access-токена без повторного ввода пароля.

В проекте время жизни access-токена задается переменной окружения \texttt{ACCESS\_TOKEN\_EXPIRE\_MINUTES}, по умолчанию --- 30 минут. Время жизни refresh-токена задается переменной \texttt{REFRESH\_TOKEN\_EXPIRE\_DAYS}, по умолчанию --- 7 дней. Это означает, что при активной работе пользователя frontend может обновлять access-токен, а после длительного бездействия потребуется повторный вход.

JWT содержит идентификатор пользователя в поле \texttt{sub}, тип токена и время истечения \texttt{exp}. Сервер подписывает токен секретным ключом. При каждом защищенном запросе FastAPI dependency \texttt{get\_current\_user} извлекает Bearer-токен, декодирует его, проверяет тип и срок действия, получает пользователя из базы данных и передает его в endpoint. Если токен недействителен, сервер возвращает ошибку 401.

\begin{figure}[H]
\centering
\begin{tikzpicture}[
    actor/.style={rectangle, draw, align=center, minimum width=35mm, minimum height=10mm},
    arrow/.style={-{Latex[length=2.5mm]}, thick},
    node distance=15mm
]
\node[actor] (client) {Frontend};
\node[actor, right=28mm of client] (api) {FastAPI};
\node[actor, right=28mm of api] (db) {PostgreSQL};

\draw[arrow] (client) -- node[above]{email, password} (api);
\draw[arrow] (api) -- node[above]{поиск user} (db);
\draw[arrow] (db) -- node[below]{password\_hash} (api);
\draw[arrow] (api) -- node[below]{access + refresh} (client);
\end{tikzpicture}
\caption{Схема входа пользователя}
\end{figure}

Такой подход является компромиссом между простотой и безопасностью. Для учебного приложения не требуется сложная корпоративная система единого входа, но базовая защита персональных данных необходима. JWT позволяет не хранить серверную сессию для каждого пользователя и хорошо подходит для REST API.

\section{Защита данных и разграничение доступа}

Защита данных в приложении строится на нескольких уровнях. Первый уровень --- хранение паролей. Пользовательский пароль никогда не сохраняется в базе данных в открытом виде. При регистрации он преобразуется в bcrypt-хэш, и именно этот хэш записывается в таблицу \texttt{users}. При входе введенный пароль снова проверяется через bcrypt, что позволяет убедиться в его правильности без восстановления исходной строки.

Второй уровень --- токеновая аутентификация. Access-токен используется для обычных защищенных запросов, а refresh-токен --- только для получения нового access-токена. Разделение этих токенов уменьшает риск долгосрочного доступа при компрометации короткоживущего access-токена. В коде backend проверяется не только подпись токена, но и его тип. Это важно, потому что refresh-токен не должен использоваться как пропуск ко всем endpoint'ам приложения.

Третий уровень --- фильтрация данных по текущему пользователю. Даже если клиент передаст идентификатор события, сервер не должен доверять ему без проверки владельца. Поэтому операции чтения, изменения и удаления должны выполняться с условием \texttt{user\_id == current\_user.id}. Такой подход защищает от ситуации, когда пользователь вручную меняет id в URL или в запросе и пытается получить доступ к чужой записи.

Четвертый уровень --- контроль данных, возвращаемых клиенту. ORM-модель пользователя содержит поле \texttt{password\_hash}, однако оно не должно попадать в JSON-ответ. Для этого используются отдельные Pydantic-схемы вывода, в которых перечислены только разрешенные поля: id, email, name, avatar и дата создания. Таким образом, внутренние служебные данные остаются на сервере.

\begin{longtable}{|p{0.25\textwidth}|p{0.33\textwidth}|p{0.3\textwidth}|}
\caption{Механизмы защиты данных}\\
\hline
\textbf{Механизм} & \textbf{Назначение} & \textbf{Реализация}\\
\hline
bcrypt & Защита пароля при хранении & Сохранение хэша вместо исходного пароля\\
\hline
JWT access-token & Доступ к защищенным endpoint'ам & Bearer-токен в HTTP-заголовке Authorization\\
\hline
JWT refresh-token & Продление сессии без повторного ввода пароля & Endpoint \texttt{/api/auth/refresh}\\
\hline
Dependency FastAPI & Получение текущего пользователя & \texttt{get\_current\_user}\\
\hline
Фильтрация по user\_id & Изоляция персональных данных & SQLAlchemy-запросы с условием владельца\\
\hline
Pydantic-схемы вывода & Исключение служебных полей из ответа & \texttt{UserOut}, \texttt{EventOut}, \texttt{NoteOut}\\
\hline
\end{longtable}

С точки зрения дальнейшего развития приложения можно усилить защиту дополнительными мерами: хранением refresh-токена в httpOnly cookie, добавлением механизма отзыва refresh-токенов, ограничением частоты запросов на вход, журналированием подозрительных попыток авторизации и более строгой настройкой CORS для production-окружения. Однако для курсового проекта реализованная схема уже покрывает базовые требования безопасности web-приложения.

\section{Реализация событий и заметок}

События и заметки являются основными пользовательскими сущностями. Их реализация строится по схожему принципу: отдельная таблица в базе данных, Pydantic-схемы для запросов и ответов, роутер FastAPI с CRUD-операциями и функции на frontend для обращения к API.

При получении списка событий клиент может передать диапазон дат. Сервер строит запрос к таблице events, фильтрует записи по текущему user\_id и указанным датам, сортирует результат и возвращает список. Фильтрация по user\_id выполняется на сервере, поэтому пользователь не может получить чужие события даже при попытке изменить параметры запроса.

Создание события включает проверку входных данных, создание ORM-объекта и сохранение в базе. При обновлении сервер сначала получает запись по id и user\_id. Это важно: если событие с таким id существует, но принадлежит другому пользователю, endpoint не должен его изменять. Удаление работает по тому же принципу.

Заметки отличаются тем, что дата у них необязательна. Это позволяет хранить как заметки, привязанные к конкретному дню, так и общие записи. На frontend они могут отображаться на отдельной странице заметок или рядом с календарем. В перспективе заметку можно преобразовывать в событие: пользователь добавляет дату и время, после чего запись становится частью расписания.

\section{Реализация клиентской части}

Клиентская часть реализована на Next.js. Страницы приложения расположены в каталоге \texttt{src/app}. Страница \texttt{auth} отвечает за вход и регистрацию, главная страница --- за отображение календаря, страница \texttt{notes} --- за работу с заметками, страница \texttt{analytics} --- за статистику и AI-виджет.

React-компоненты позволяют разделить интерфейс на самостоятельные части. Например, Header отображает верхнюю панель, Calendar отвечает за календарную сетку, EventModal --- за создание или редактирование события, AIWidget --- за запрос и отображение нейросетевого анализа. Такой подход уменьшает связанность кода: изменение виджета аналитики не требует переписывать страницу авторизации.

Для обмена с сервером используется модуль \texttt{src/lib/api.ts}. Он содержит функции для регистрации, входа, получения событий, создания заметок, получения аналитики и AI-анализа. В этом же модуле реализована обработка access- и refresh-токенов. Если сервер возвращает ошибку авторизации, клиент пытается обновить access-токен; если обновление невозможно, пользователь перенаправляется на страницу входа.

TypeScript-типы описывают структуру данных на клиенте. Например, тип \texttt{CalendarEvent} содержит id, userId, title, date, startTime, endTime, type и другие поля. Наличие типов уменьшает риск ошибки при переименовании полей или изменении структуры ответа backend.

Визуальное оформление построено на Tailwind CSS. Этот инструмент особенно удобен для учебного проекта, потому что позволяет быстро создавать адаптивные интерфейсы без сложной CSS-архитектуры. Например, карточки событий, кнопки, сетки и панели можно стилизовать прямо в JSX-разметке с помощью утилитарных классов.

\section{Реализация клиентского API-слоя}

Отдельное место в клиентской части занимает слой взаимодействия с backend. Его целесообразно выделять в отдельный модуль, потому что запросы к серверу используются на разных страницах: авторизация нужна на странице входа, события --- на календарной странице, заметки --- на странице заметок, аналитика и AI --- на странице аналитики. Если каждая страница будет самостоятельно формировать \texttt{fetch}-запросы, код быстро станет дублироваться и усложнит обработку ошибок.

В проекте функции обращения к серверу сгруппированы в модуле \texttt{src/lib/api.ts}. Такой модуль выполняет несколько задач. Во-первых, он хранит базовый адрес backend. Во-вторых, он добавляет access-токен в заголовок \texttt{Authorization}. В-третьих, он преобразует данные из JSON в TypeScript-объекты, с которыми дальше работает интерфейс. В-четвертых, он обрабатывает ошибки авторизации и может выполнить запрос на обновление токена.

Типовая цепочка запроса выглядит следующим образом. Компонент страницы вызывает функцию, например \texttt{getEvents}. Эта функция формирует URL с параметрами периода, добавляет заголовки и отправляет запрос. Если ответ успешен, JSON преобразуется в массив событий. Если сервер возвращает ошибку 401, клиент пытается обновить access-токен через refresh-токен. Если обновление успешно, исходный запрос повторяется. Если refresh-токен также недействителен, пользователь перенаправляется на страницу авторизации.

\begin{figure}[H]
\centering
\begin{tikzpicture}[
    box/.style={rectangle, draw, rounded corners, align=center, minimum width=35mm, minimum height=10mm},
    arrow/.style={-{Latex[length=2.5mm]}, thick},
    node distance=13mm
]
\node[box] (page) {React-страница};
\node[box, below=of page] (api) {\texttt{src/lib/api.ts}};
\node[box, below=of api] (backend) {FastAPI backend};
\node[box, right=25mm of api] (storage) {access/refresh\\token};

\draw[arrow] (page) -- node[right]{вызов функции} (api);
\draw[arrow] (api) -- node[right]{HTTP-запрос} (backend);
\draw[arrow] (api) -- node[above]{чтение/запись} (storage);
\draw[arrow] (backend) -- ++(-40mm,0) |- node[left]{JSON-ответ} (page);
\end{tikzpicture}
\caption{Слой взаимодействия frontend с backend}
\end{figure}

Такой слой делает клиентскую часть более устойчивой. Например, при изменении адреса backend или схемы обновления токена достаточно изменить один модуль, а не все страницы приложения. Кроме того, это упрощает тестирование: можно отдельно проверить функции API и отдельно визуальные компоненты.

\section{Реализация основных пользовательских сценариев}

Реализация приложения проверялась через пользовательские сценарии, так как именно они показывают, насколько связно работают отдельные модули. Сценарий регистрации начинается на frontend: пользователь вводит имя, email и пароль. После отправки формы клиент вызывает endpoint регистрации. Backend проверяет email, хэширует пароль и создает пользователя. После успешного ответа пользователь может перейти ко входу или сразу получить токены в зависимости от выбранной логики приложения.

Сценарий входа похож, но вместо создания пользователя выполняется проверка учетных данных. После получения access- и refresh-токенов клиент сохраняет их и использует access-токен для дальнейших запросов. Затем выполняется запрос \texttt{/api/auth/me}, чтобы получить имя и email текущего пользователя для верхней панели.

Сценарий создания события включает выбор даты, заполнение названия, типа, времени и необязательных полей. Frontend отправляет данные на endpoint \texttt{/api/events}. Backend валидирует их через Pydantic, добавляет \texttt{user\_id} текущего пользователя и сохраняет запись. После ответа клиент обновляет календарь, чтобы пользователь сразу увидел новое событие.

Сценарий аналитики начинается с выбора периода или использования периода по умолчанию. Клиент запрашивает агрегированные данные у backend. Сервер выбирает события и заметки текущего пользователя, вычисляет показатели и возвращает результат. На странице аналитики эти данные отображаются в виде карточек и списков. При запросе AI-анализа frontend обращается к отдельному endpoint'у, потому что генерация текста может занимать больше времени, чем обычный SQL-запрос.

\begin{longtable}{|p{0.22\textwidth}|p{0.34\textwidth}|p{0.32\textwidth}|}
\caption{Пользовательские сценарии и задействованные модули}\\
\hline
\textbf{Сценарий} & \textbf{Действия пользователя} & \textbf{Модули приложения}\\
\hline
Регистрация & Ввод имени, email и пароля & Auth page, \texttt{/api/auth/register}, модель \texttt{User}\\
\hline
Вход & Ввод email и пароля, получение токенов & Auth page, \texttt{/api/auth/login}, JWT-сервис\\
\hline
Создание события & Заполнение формы события & Calendar UI, \texttt{/api/events}, модель \texttt{Event}\\
\hline
Создание заметки & Ввод заголовка и текста & Notes page, \texttt{/api/notes}, модель \texttt{Note}\\
\hline
Просмотр аналитики & Переход на страницу аналитики & Analytics page, \texttt{/api/analytics}\\
\hline
AI-анализ & Нажатие кнопки анализа & AIWidget, \texttt{/api/ai/analysis}, Ollama\\
\hline
\end{longtable}

\section{Реализация аналитики}

Аналитический модуль решает задачу обобщения календарных данных. Он предоставляет пользователю числовые показатели: общее количество событий, распределение по типам, активность по дням и количество заметок. Эти данные помогают оценить загрузку и выявить периоды с повышенной активностью.

На backend аналитика реализована отдельным роутером. Endpoint принимает диапазон дат и текущего пользователя, выбирает события за период, группирует их по типам и формирует список ежедневной активности. Количество заметок считается отдельно. Результат возвращается в виде Pydantic-схемы \texttt{AnalyticsOut}.

На frontend страница аналитики отображает показатели в карточках и графиках. Пользователь видит не просто сырой список событий, а агрегированную картину. Это отличает приложение от простого календаря: система начинает выполнять аналитическую функцию и помогает пользователю понять, как распределяется его время.

Аналитика также служит подготовительным уровнем для AI-модуля. Числовая статистика полезна, но текстовый анализ может дать более понятные рекомендации. Например, если на неделе много задач одного типа, нейросеть может предложить сгруппировать их или выделить отдельные временные блоки.

\section{Интеграция локальной нейросети}

Одной из особенностей проекта является интеграция нейросетевого анализа. В отличие от варианта, где данные отправляются во внешний облачный сервис, в проекте используется локальная среда Ollama. Backend обращается к ней через OpenAI-compatible API. Это значит, что код клиента использует библиотеку \texttt{AsyncOpenAI}, но адрес API указывает на локальный сервер Ollama.

Такой подход имеет несколько преимуществ. Во-первых, пользовательские события и заметки не покидают локальное окружение, что лучше соответствует идее прозрачной обработки данных. Во-вторых, приложение не требует платного облачного API для демонстрации функции. В-третьих, можно менять модель, например использовать qwen2.5:1.5b или другую локально установленную модель.

Сервис AI-анализа получает список событий и заметок. Затем он преобразует их в текстовый контекст: для каждого события указывается дата, время, название и тип; для каждой заметки --- заголовок и текст. После этого формируется промпт, в котором модели задается роль персонального AI-ассистента по продуктивности. Модель должна кратко проанализировать нагрузку пользователя, выделить основные темы и дать практические советы.

\begin{figure}[H]
\centering
\begin{tikzpicture}[
    box/.style={rectangle, draw, rounded corners, align=center, minimum width=38mm, minimum height=11mm},
    arrow/.style={-{Latex[length=2.5mm]}, thick},
    node distance=13mm
]
\node[box] (ui) {Страница аналитики};
\node[box, below=of ui] (router) {FastAPI router /api/ai};
\node[box, below=of router] (service) {AI service\\формирование промпта};
\node[box, below=of service] (ollama) {Ollama\\локальная LLM};
\node[box, right=25mm of service] (db) {PostgreSQL\\events, notes};

\draw[arrow] (ui) -- node[right]{запрос анализа} (router);
\draw[arrow] (router) -- (service);
\draw[arrow] (router) -- node[above]{данные пользователя} (db);
\draw[arrow] (service) -- node[right]{prompt} (ollama);
\draw[arrow] (ollama) -- ++(-35mm,0) |- node[left]{текстовый отчет} (ui);
\end{tikzpicture}
\caption{Взаимодействие с локальной нейросетью}
\end{figure}

При недоступности Ollama сервис возвращает понятное сообщение об ошибке. Это важно для пользовательского опыта: вместо технического stack trace пользователь получает объяснение, что локальная нейросеть не запущена.

\section{Формирование промпта и обработка ответа нейросети}

Нейросетевой модуль не получает прямой доступ к базе данных. Данные выбираются серверной частью приложения, после чего передаются в сервис формирования промпта. Это решение важно с точки зрения контроля: backend сам определяет, какие сведения пользователя можно отправить в модель, в каком объеме и в каком виде. Нейросеть не выполняет SQL-запросы и не может самостоятельно обратиться к данным другого пользователя.

Промпт состоит из нескольких частей. Первая часть задает роль модели: персональный ассистент по продуктивности в приложении ``Умный календарь''. Вторая часть содержит события пользователя за выбранный период. Для каждого события указываются дата, время, название и тип. Третья часть содержит заметки пользователя. Четвертая часть задает инструкцию к формату ответа: кратко проанализировать нагрузку, выделить основные темы и дать практические рекомендации на русском языке.

Такой подход делает поведение модели более предсказуемым. Если передать модели только сырой JSON, она может ответить слишком технически или уделить внимание не тем полям. Если же сформировать текстовый контекст с понятной ролью и инструкциями, вероятность полезного ответа повышается. При этом промпт остается достаточно коротким, что важно при использовании локальной модели с ограниченными ресурсами.

\begin{longtable}{|p{0.24\textwidth}|p{0.3\textwidth}|p{0.34\textwidth}|}
\caption{Структура промпта для AI-анализа}\\
\hline
\textbf{Часть промпта} & \textbf{Содержимое} & \textbf{Назначение}\\
\hline
Роль & Описание модели как персонального ассистента & Задает стиль и предметную область ответа\\
\hline
События & Даты, время, названия и типы событий & Позволяет оценить календарную нагрузку\\
\hline
Заметки & Заголовки и тексты заметок & Позволяет выявить темы, не отраженные в расписании\\
\hline
Инструкция & Требования к языку, краткости и советам & Ограничивает форму итогового ответа\\
\hline
\end{longtable}

Ответ нейросети возвращается пользователю как текстовый отчет. В текущей реализации он не изменяет календарь автоматически. Это осознанное ограничение: модель может ошибаться, поэтому на данном этапе она выполняет рекомендательную функцию. Пользователь сам принимает решение, стоит ли переносить событие, добавлять напоминание или менять распределение задач.

В дальнейшем AI-модуль можно развить в интерактивного помощника. Например, пользователь сможет задавать вопросы: ``Какие дни у меня самые загруженные?'', ``Когда лучше поставить подготовку к экзамену?'', ``Какие задачи можно перенести?''. Для такого режима потребуется не только генерация текста, но и безопасный механизм подтверждения действий перед изменением календаря.

\section{Контейнеризация и запуск}

Для упрощения запуска серверной части и базы данных используется Docker Compose. В конфигурации описан сервис PostgreSQL и сервис backend. Backend зависит от готовности базы данных, что позволяет избежать ситуации, когда приложение стартует раньше, чем СУБД готова принимать подключения.

Контейнеризация полезна не только для разработки, но и для демонстрации курсового проекта. Она уменьшает количество действий, которые нужно выполнить для запуска: не требуется отдельно устанавливать и настраивать PostgreSQL, создавать базу данных и вручную прописывать параметры подключения. Все основные настройки передаются через переменные окружения.

Frontend запускается отдельно командой \texttt{npm run dev}. Такой вариант удобен во время разработки, так как Next.js предоставляет быструю перезагрузку страницы при изменении компонентов. В дальнейшем frontend также может быть добавлен в Docker Compose как отдельный сервис.

\section{Тестирование и проверка работоспособности}

Проверка работоспособности приложения проводилась по основным пользовательским сценариям. Сначала проверялась регистрация нового пользователя и вход в систему. После успешной авторизации проверялась загрузка данных текущего пользователя, создание событий, отображение событий в календаре, редактирование и удаление.

Далее проверялся модуль заметок. Создавались заметки с датой и без даты, после чего проверялось их отображение в интерфейсе и сохранение в базе данных. Отдельное внимание уделялось разграничению данных: запросы выполнялись от имени текущего пользователя, а серверная часть фильтровала записи по идентификатору пользователя из JWT.

После этого проверялась аналитика. Для тестового периода создавались события разных типов. Страница аналитики должна была корректно показать общее количество событий, распределение по типам и ежедневную активность. Затем проверялся AI-виджет: при запущенной локальной модели он возвращал текстовый анализ, при недоступности Ollama --- сообщение о невозможности подключения.

Наконец, проверялся запуск инфраструктуры. Backend и PostgreSQL поднимались через Docker Compose, а frontend запускался в режиме разработки. В результате были подтверждены основные сценарии: вход, регистрация, работа с календарем, заметками, аналитикой и AI-анализом.

\begin{longtable}{|p{0.24\textwidth}|p{0.3\textwidth}|p{0.34\textwidth}|}
\caption{Матрица функционального тестирования}\\
\hline
\textbf{Проверка} & \textbf{Действия} & \textbf{Ожидаемый результат}\\
\hline
Регистрация & Ввести новый email, имя и пароль & Пользователь создан, пароль сохранен как хэш, доступен вход\\
\hline
Повторная регистрация & Ввести уже существующий email & Сервер возвращает сообщение о занятости email\\
\hline
Вход & Ввести корректные учетные данные & Клиент получает access- и refresh-токены\\
\hline
Неверный пароль & Ввести неправильный пароль & Сервер возвращает ошибку авторизации\\
\hline
Обновление токена & Дождаться истечения access-токена и выполнить запрос & Клиент получает новый access-токен через refresh-токен\\
\hline
Создание события & Заполнить форму события и сохранить & Событие появляется в календаре и сохраняется в базе данных\\
\hline
Редактирование события & Изменить название или дату события & В календаре отображаются обновленные данные\\
\hline
Удаление события & Удалить выбранное событие & Событие исчезает из интерфейса и базы данных\\
\hline
Создание заметки & Добавить заметку с датой или без даты & Заметка отображается в списке заметок\\
\hline
Изоляция данных & Выполнить запрос от имени другого пользователя & Чужие события и заметки не возвращаются\\
\hline
Аналитика & Создать события разных типов и открыть аналитику & Показаны корректные агрегированные показатели\\
\hline
AI-анализ & Запустить Ollama и запросить анализ & Возвращается текстовый отчет по событиям и заметкам\\
\hline
Недоступность AI & Остановить Ollama и запросить анализ & Пользователь получает понятное сообщение об ошибке\\
\hline
\end{longtable}

Помимо функциональной проверки, важна проверка устойчивости интерфейса. Для этого следует открывать страницы приложения после перезагрузки браузера, проверять поведение при отсутствии токенов, при истекшем refresh-токене и при временно недоступном backend. Такие ситуации часто встречаются в реальной эксплуатации и позволяют обнаружить ошибки, которые не видны при обычном успешном сценарии.

Также была проведена проверка соответствия структуры проекта целям сопровождения. Серверные роутеры отделены друг от друга, схемы данных вынесены в общий файл, клиентские функции API сгруппированы отдельно от визуальных компонентов. Это позволяет развивать приложение дальше: например, добавить повторяющиеся события без переписывания авторизации или улучшить страницу аналитики без изменения модели пользователя.

\section{Сравнение реализованного результата с требованиями}

Реализованное приложение соответствует основным функциональным требованиям, сформулированным во второй главе. Поддерживается регистрация и вход, реализована работа с JWT, события и заметки сохраняются в PostgreSQL, клиентская часть отображает данные через web-интерфейс, аналитика строится на основе реальных записей пользователя.

Некоторые возможности заложены архитектурно, но могут быть развиты в дальнейшем. Например, таблицы шаблонов и уведомлений уже присутствуют в модели данных, однако в будущем можно расширить их пользовательские сценарии: добавить полноценный экран управления шаблонами, систему напоминаний по расписанию, отправку email-уведомлений или push-уведомлений.

Особое значение имеет AI-модуль. Он не является обязательным для базового календаря, но делает приложение именно ``умным'': система не только хранит информацию, но и помогает ее осмыслить. Такой подход соответствует современному направлению развития персональных цифровых ассистентов.

\section{Выводы по главе 3}

В третьей главе была описана практическая реализация web-приложения ``Умный календарь''. Рассмотрена структура проекта, серверная часть на FastAPI, Pydantic-схемы, SQLAlchemy-модели, JWT-аутентификация, клиентская часть на Next.js и TypeScript, аналитика и интеграция локальной нейросети. Показано, что приложение реализует основные сценарии работы пользователя и может служить основой для дальнейшего развития.

\chapter*{ЗАКЛЮЧЕНИЕ}
\addcontentsline{toc}{chapter}{ЗАКЛЮЧЕНИЕ}

В ходе выполнения курсового проекта было реализовано web-приложение ``Умный календарь'', предназначенное для ведения личного расписания, управления событиями и заметками, анализа пользовательской активности и получения интеллектуальных рекомендаций. В отличие от предыдущего этапа работы, где основной акцент был сделан на обзоре предметной области и формировании концепции, данная работа была направлена на практическую реализацию приложения.

В первой главе была рассмотрена предметная область календарного планирования и проведен анализ существующих решений. Было показано, что Google Calendar, Doodle и Tweek обладают полезными идеями, но не полностью соответствуют цели создания простого независимого приложения для личного планирования. На основе анализа были выявлены основные проблемы: перегруженность интерфейса, фрагментация календаря и заметок, зависимость от внешних аккаунтов и недостаток персональной аналитики.

Во второй главе были сформулированы функциональные и нефункциональные требования к приложению. Также была спроектирована архитектура системы: frontend на Next.js, backend на FastAPI, база данных PostgreSQL и локальная нейросеть Ollama как внешний интеллектуальный компонент. Были определены основные сущности базы данных и ключевые endpoint'ы REST API.

В третьей главе была описана реализация приложения. Серверная часть разработана на Python с использованием FastAPI, SQLAlchemy и Pydantic. Реализованы модели пользователей, событий, заметок, шаблонов и уведомлений. Для защиты данных используется JWT-аутентификация с access- и refresh-токенами, а пароли хранятся в виде bcrypt-хэшей. Клиентская часть разработана на Next.js, React и TypeScript, а визуальное оформление выполнено с помощью Tailwind CSS. Реализованы страницы авторизации, календаря, заметок и аналитики. Отдельно разработан AI-модуль, который формирует текстовый анализ активности пользователя на основе его событий и заметок.

Поставленная цель курсового проекта достигнута: создано работоспособное web-приложение, объединяющее календарь, заметки, аналитику и интеллектуальные рекомендации. Практическая значимость работы заключается в том, что разработанное решение можно использовать как основу персонального планировщика или развивать дальше в рамках дипломной работы. Архитектура приложения допускает добавление новых возможностей: группового планирования, уведомлений, импорта и экспорта iCal, мобильной версии, расширенной аналитики и интеграций с внешними календарями.

В качестве дальнейших направлений развития можно выделить реализацию полноценного механизма напоминаний, настройку повторяющихся событий, поддержку нескольких календарей, улучшение UI для мобильных устройств, добавление ролей пользователей, экспорт данных, а также расширение AI-модуля до режима диалогового помощника, который не только анализирует расписание, но и предлагает конкретные изменения в календаре.

\begin{thebibliography}{99}
\addcontentsline{toc}{chapter}{СПИСОК ИСПОЛЬЗОВАННЫХ ИСТОЧНИКОВ}

\bibitem{fielding}
Fielding, R. T. Architectural Styles and the Design of Network-based Software Architectures : doctoral dissertation / R. T. Fielding. -- University of California, Irvine, 2000. -- 180 p.

\bibitem{fastapi}
FastAPI Documentation [Электронный ресурс]. -- Режим доступа: \url{https://fastapi.tiangolo.com/}. -- Дата доступа: 15.05.2026.

\bibitem{nextjs}
Next.js Documentation [Электронный ресурс]. -- Режим доступа: \url{https://nextjs.org/docs}. -- Дата доступа: 15.05.2026.

\bibitem{react}
React Documentation [Электронный ресурс]. -- Режим доступа: \url{https://react.dev/}. -- Дата доступа: 15.05.2026.

\bibitem{typescript}
TypeScript Documentation [Электронный ресурс]. -- Режим доступа: \url{https://www.typescriptlang.org/docs/}. -- Дата доступа: 15.05.2026.

\bibitem{sqlalchemy}
SQLAlchemy 2.0 Documentation [Электронный ресурс]. -- Режим доступа: \url{https://docs.sqlalchemy.org/en/20/}. -- Дата доступа: 15.05.2026.

\bibitem{postgresql}
PostgreSQL Documentation [Электронный ресурс]. -- Режим доступа: \url{https://www.postgresql.org/docs/}. -- Дата доступа: 15.05.2026.

\bibitem{pydantic}
Pydantic Documentation [Электронный ресурс]. -- Режим доступа: \url{https://docs.pydantic.dev/}. -- Дата доступа: 15.05.2026.

\bibitem{jwt}
Jones, M. JSON Web Token (JWT). RFC 7519 [Электронный ресурс] / M. Jones, J. Bradley, N. Sakimura. -- Режим доступа: \url{https://www.rfc-editor.org/rfc/rfc7519}. -- Дата доступа: 15.05.2026.

\bibitem{tailwind}
Tailwind CSS Documentation [Электронный ресурс]. -- Режим доступа: \url{https://tailwindcss.com/docs}. -- Дата доступа: 15.05.2026.

\bibitem{docker}
Docker Documentation [Электронный ресурс]. -- Режим доступа: \url{https://docs.docker.com/}. -- Дата доступа: 15.05.2026.

\bibitem{openai}
OpenAI API Documentation [Электронный ресурс]. -- Режим доступа: \url{https://platform.openai.com/docs}. -- Дата доступа: 15.05.2026.

\bibitem{ollama}
Ollama Documentation [Электронный ресурс]. -- Режим доступа: \url{https://ollama.com/}. -- Дата доступа: 15.05.2026.

\bibitem{norman}
Норман, Д. Дизайн привычных вещей / Д. Норман ; пер. с англ. -- М. : Манн, Иванов и Фербер, 2018. -- 384 с.

\bibitem{fowler}
Фаулер, М. Архитектура корпоративных программных приложений / М. Фаулер ; пер. с англ. -- СПб. : Питер, 2020. -- 544 с.

\end{thebibliography}

\appendix

\chapter*{ПРИЛОЖЕНИЕ А\\ФРАГМЕНТЫ СТРУКТУРЫ REST API}
\addcontentsline{toc}{chapter}{ПРИЛОЖЕНИЕ А. ФРАГМЕНТЫ СТРУКТУРЫ REST API}

\begin{Verbatim}[fontsize=\small]
POST   /api/auth/register
POST   /api/auth/login
POST   /api/auth/refresh
GET    /api/auth/me

GET    /api/events
POST   /api/events
GET    /api/events/{id}
PUT    /api/events/{id}
DELETE /api/events/{id}

GET    /api/notes
POST   /api/notes
PUT    /api/notes/{id}
DELETE /api/notes/{id}

GET    /api/analytics
GET    /api/ai/analysis
\end{Verbatim}

Приведенная структура демонстрирует REST-ориентированный характер API. Endpoint'ы сгруппированы по предметным областям, а операции над сущностями используют стандартные HTTP-методы. Такой подход облегчает сопровождение приложения и делает API понятным для возможных внешних клиентов.

\chapter*{ПРИЛОЖЕНИЕ Б\\ОПИСАНИЕ ОСНОВНЫХ СХЕМ ДАННЫХ}
\addcontentsline{toc}{chapter}{ПРИЛОЖЕНИЕ Б. ОПИСАНИЕ ОСНОВНЫХ СХЕМ ДАННЫХ}

\begin{longtable}{|p{0.22\textwidth}|p{0.68\textwidth}|}
\hline
\textbf{Схема} & \textbf{Назначение}\\
\hline
UserCreate & Данные для регистрации пользователя: email, имя, пароль.\\
\hline
UserLogin & Данные для входа пользователя: email и пароль.\\
\hline
Token & Ответ сервера после входа: access\_token, refresh\_token, token\_type.\\
\hline
EventCreate & Данные для создания события: название, дата, время, тип, описание, ссылка.\\
\hline
EventUpdate & Частичное обновление события; все поля являются необязательными.\\
\hline
EventOut & Ответ API с полным описанием события, включая id и дату создания.\\
\hline
NoteCreate & Данные для создания заметки: заголовок, текст, необязательная дата.\\
\hline
AnalyticsOut & Статистика пользователя: общее число событий, распределение по типам, активность по дням, число заметок.\\
\hline
\end{longtable}

\chapter*{ПРИЛОЖЕНИЕ В\\ПРИМЕР ПРОМПТА ДЛЯ НЕЙРОСЕТЕВОГО АНАЛИЗА}
\addcontentsline{toc}{chapter}{ПРИЛОЖЕНИЕ В. ПРИМЕР ПРОМПТА ДЛЯ НЕЙРОСЕТЕВОГО АНАЛИЗА}

\begin{Verbatim}[fontsize=\small]
Ты -- персональный AI-ассистент в приложении "Умный Календарь".
Твоя задача: проанализировать данные пользователя за неделю
и дать краткий, но полезный отчет о продуктивности.

СПИСОК СОБЫТИЙ И ЗАДАЧ:
- 2026-05-12 10:00: Подготовить доклад (task)
- 2026-05-13 14:00: Встреча по проекту (meeting)

ТЕКСТЫ ЗАМЕТОК:
- Курсовая: проверить список литературы
- Идея: добавить напоминания

Инструкция:
1. Проанализируй нагрузку пользователя.
2. Выдели главные темы из заметок и задач.
3. Дай 2 коротких практических совета.
4. Отвечай на русском языке.
\end{Verbatim}

Промпт формируется на сервере динамически. События и заметки подставляются из базы данных текущего пользователя, поэтому отчет отражает реальную активность, а не заранее подготовленный шаблон.

\chapter*{ПРИЛОЖЕНИЕ Г\\АЛГОРИТМ ОБНОВЛЕНИЯ ACCESS-ТОКЕНА}
\addcontentsline{toc}{chapter}{ПРИЛОЖЕНИЕ Г. АЛГОРИТМ ОБНОВЛЕНИЯ ACCESS-ТОКЕНА}

\begin{enumerate}
    \item Пользователь выполняет вход и получает два токена: \texttt{access\_token} и \texttt{refresh\_token}.
    \item Клиент сохраняет токены и добавляет \texttt{access\_token} в заголовок \texttt{Authorization} при обращении к защищенным endpoint'ам.
    \item Если access-токен действителен, backend выполняет запрос и возвращает данные.
    \item Если срок действия access-токена истек, backend возвращает ошибку авторизации.
    \item Клиент отправляет refresh-токен на endpoint \texttt{/api/auth/refresh}.
    \item Backend проверяет подпись refresh-токена, его тип и срок действия, после чего находит пользователя в базе данных.
    \item Если refresh-токен действителен, backend формирует новый access-токен и возвращает его клиенту.
    \item Клиент повторяет исходный запрос уже с новым access-токеном.
    \item Если refresh-токен недействителен или истек, пользователь перенаправляется на страницу входа.
\end{enumerate}

\begin{figure}[H]
\centering
\begin{tikzpicture}[
    box/.style={rectangle, draw, rounded corners, align=center, minimum width=40mm, minimum height=10mm},
    arrow/.style={-{Latex[length=2.5mm]}, thick},
    node distance=12mm
]
\node[box] (req) {Запрос к API};
\node[box, below=of req] (check) {Проверка access-токена};
\node[box, below left=14mm and 8mm of check] (ok) {Токен действителен\\данные возвращены};
\node[box, below right=14mm and 8mm of check] (bad) {Токен истек\\ошибка 401};
\node[box, below=of bad] (refresh) {Запрос refresh};
\node[box, below=of refresh] (new) {Новый access-токен};

\draw[arrow] (req) -- (check);
\draw[arrow] (check) -- node[left]{да} (ok);
\draw[arrow] (check) -- node[right]{нет} (bad);
\draw[arrow] (bad) -- (refresh);
\draw[arrow] (refresh) -- (new);
\draw[arrow] (new) -| (req);
\end{tikzpicture}
\caption{Обновление access-токена через refresh-токен}
\end{figure}

\chapter*{ПРИЛОЖЕНИЕ Д\\ИНСТРУКЦИЯ ПО ЗАПУСКУ ПРИЛОЖЕНИЯ}
\addcontentsline{toc}{chapter}{ПРИЛОЖЕНИЕ Д. ИНСТРУКЦИЯ ПО ЗАПУСКУ ПРИЛОЖЕНИЯ}

Для запуска приложения в локальном окружении необходимо подготовить backend, базу данных, frontend и, при необходимости, локальную нейросеть.

\begin{enumerate}
    \item Перейти в каталог проекта \texttt{Kursovoi\_Project}.
    \item Запустить серверную часть и PostgreSQL через Docker Compose:
\begin{Verbatim}[fontsize=\small]
docker compose up -d
\end{Verbatim}
    \item Перейти в каталог \texttt{frontend} и установить зависимости, если они еще не установлены:
\begin{Verbatim}[fontsize=\small]
npm install
\end{Verbatim}
    \item Запустить frontend в режиме разработки:
\begin{Verbatim}[fontsize=\small]
npm run dev
\end{Verbatim}
    \item Открыть приложение в браузере по адресу \texttt{http://localhost:3000}.
    \item Для работы AI-анализа запустить Ollama и убедиться, что локальная модель доступна по адресу \texttt{http://localhost:11434}.
\end{enumerate}

Если при входе или регистрации появляется ошибка сетевого запроса, в первую очередь нужно проверить, запущен ли backend, совпадает ли адрес API в настройках frontend и открыт ли нужный порт. Если ошибка возникает только на странице AI-анализа, следует проверить состояние Ollama и наличие выбранной локальной модели.

\chapter*{ПРИЛОЖЕНИЕ Е\\ВОЗМОЖНЫЕ НАПРАВЛЕНИЯ РАЗВИТИЯ}
\addcontentsline{toc}{chapter}{ПРИЛОЖЕНИЕ Е. ВОЗМОЖНЫЕ НАПРАВЛЕНИЯ РАЗВИТИЯ}

\begin{longtable}{|p{0.28\textwidth}|p{0.58\textwidth}|}
\hline
\textbf{Направление} & \textbf{Описание}\\
\hline
Повторяющиеся события & Добавление правил повторения для занятий, лабораторных работ, еженедельных встреч и регулярных задач.\\
\hline
Напоминания & Реализация фоновой задачи, которая будет создавать и отправлять уведомления перед началом события.\\
\hline
Импорт и экспорт iCal & Возможность переносить события между ``Умным календарем'' и внешними календарными сервисами.\\
\hline
Мобильная адаптация & Улучшение интерфейса для небольших экранов и сценариев быстрого добавления события.\\
\hline
Диалоговый AI-помощник & Возможность задавать вопросы по расписанию и получать предложения по оптимизации календаря.\\
\hline
Групповое планирование & Создание общих событий и подбор свободного времени для нескольких пользователей.\\
\hline
Расширенная аналитика & Анализ недельных и месячных трендов, перегруженных дней, баланса учебных и личных задач.\\
\hline
\end{longtable}

\end{document}
